\section{Tabelle}
\begin{table}[H]
\centering
\caption{Portfolio representations generated by different kernels}
\label{tab:kernel_policy_classes}
\small
\renewcommand{\arraystretch}{1.25}

\begin{tabular}{
p{0.21\textwidth}
p{0.34\textwidth}
p{0.35\textwidth}
}
\toprule
Representation & Kernel & Portfolio policy\\
\midrule

Common exposure
&
$K(z,z')=1$
&
$W(z)$ is constant
\\[4pt]

Linear
&
$K(z,z')=1+z^\prime z'$
&
$W(z)$ is linear in characteristics
\\[4pt]

Finite features
&
$K(z,z')=\Phi(z)^\prime\Phi(z')$
&
$W(z)$ lies in the span of the chosen features
\\[4pt]

Polynomial
&
$K(z,z')=(1+z^\prime z')^q$
&
$W(z)$ includes polynomial terms and interactions
\\[6pt]

Gaussian
&
$\displaystyle
K(z,z')
=
\exp\!\left(
-\frac{\|z-z'\|^2}{2\ell^2}
\right)$
&
Flexible nonlinear $W(z)$
\\[8pt]

Mat\'ern
&
$\displaystyle
\begin{aligned}
K_s(z,z')
&=
c_{s,D}\,
\rho_s(z,z')^{\,s-D/2}
\\[-1pt]
&\quad\times
\mathcal K_{s-D/2}
\!\bigl(\rho_s(z,z')\bigr)
\end{aligned}$
&
Flexible nonlinear $W(z)$ with smoothness controlled by $s$
\\[10pt]

Neural tangent
&
$\displaystyle
\begin{aligned}
K_{\mathrm{NTK}}(z,z')
&=
\frac{\|z\|\,\|z'\|}{2\pi}
\Bigl[\sin\theta
\\[-1pt]
&\qquad
+2(\pi-\theta)\cos\theta
\Bigr]
\end{aligned}$
&
Infinite-width one-hidden-layer ReLU policy
\\

\bottomrule
\end{tabular}

\begin{minipage}{0.97\textwidth}
\footnotesize
\emph{Notes:}
Here $\ell>0$ denotes the kernel length scale. For the Mat\'ern kernel,
\[
\rho_s(z,z')
=
\sqrt{2s-D}\,
\frac{\|z-z'\|}{\ell},
\qquad
c_{s,D}
=
\frac{2^{1-(s-D/2)}}{\Gamma(s-D/2)},
\qquad
s>\frac{D}{2},
\]
and $\mathcal K_\alpha$ denotes the modified Bessel function of the second kind.
Under standard regularity conditions, the Mat\'ern RKHS is equivalent to a
Sobolev space of smoothness $s$.

For the neural tangent kernel,
\[
\theta
=
\arccos\!\left(
\frac{z^\prime z'}
{\|z\|\,\|z'\|}
\right),
\]
and the displayed expression corresponds to a one-hidden-layer ReLU network
without bias in the infinite-width NTK limit, up to the conventional overall
normalization.

The kernel determines the class of characteristic-to-exposure mappings
available to the portfolio manager. Portfolio returns use the common
cross-sectional normalization in Equation~\eqref{eq:portfolio_return}.
The common-exposure specification is an equal-exposure characteristic rule,
not an unrestricted vector of static stock weights.
\end{minipage}
\end{table}


\Needspace{15\baselineskip}
Table~\ref{tab:special-rates} relates the learning exponent to familiar
spectral benchmarks. These are conditions on the managed-portfolio operator,
not conclusions implied by the kernel name alone.

\begin{table}[H]
\centering
\caption{Benchmark portfolio-learning rates}
\label{tab:special-rates}
\small
\begin{tabular}{lcc}
\toprule
Portfolio class
&
Managed-portfolio spectrum
&
Sharpe learning rate
\\
\midrule
Finite dimensional
&
Fixed finite rank
&
$T^{-1}$
\\[5pt]

Gaussian benchmark
&
Fast decay
&
Nearly $T^{-1}$
\\[5pt]

Sobolev--Mat\'ern benchmark
&
$b=2\widetilde s/D$
&
$T^{-2\widetilde s/(2\widetilde s+D)}$
\\
\bottomrule
\end{tabular}

\par\smallskip
\begin{minipage}{0.97\textwidth}
\footnotesize
\emph{Notes:} Sharpe rates use Condition~\ref{ec:sharpe}.
The finite-rank rate uses a fixed identifiable payoff space and root-$T$
moment estimation. The fast-decay row is a separate i.i.d. or sufficiently
fast-mixing benchmark, with possible logarithmic factors and an admissible
regularization sequence; it does not follow by setting $b=\infty$ in
Theorem~\ref{thm:rate}. The Sobolev--Mat\'ern row assumes the displayed
\emph{managed-payoff} spectral law and $\gamma>b$.
\end{minipage}
\end{table}
